% =========================================================
% Axiom-Derived Set Operations
% =========================================================

\subsection{Axiom-Derived Set Operations}

% ---------------------------------------------------------
% TOOLKIT
% ---------------------------------------------------------

\begin{lratopiccard}{Axiom-Derived Operations --- Construction Signature}
\begin{lrasignaturetable}
\LraSignatureRow{Native data}{Objects are sets; the only primitive relation is membership \(x\in A\).}
\LraSignatureRow{Axiom-licensed outputs}{\(\varnothing\), \(\{a,b\}\), \(\bigcup A\), \(\mathcal P(A)\), and separated subsets \(\{x\in A\mid \varphi(x)\}\).}
\LraSignatureRow{Derived outputs}{\(A\cup B\), \(A\cap B\), \(A\setminus B\), \(A\triangle B\), \(A^c\), and \(A\times B\).}
\LraSignatureRow{Licensing proof}{State the intended membership profile; prove existence; prove uniqueness by Extensionality; then introduce notation.}
\LraSignatureRow{Use rule}{Every later calculation begins by rewriting \(x\in\) the constructed set into its membership condition.}
\end{lrasignaturetable}
\end{lratopiccard}

% ---------------------------------------------------------
% Well-foundedness for definitional extensions
% ---------------------------------------------------------
\begin{theorembox}{Theorem (Existence and Uniqueness of the Empty Set)}
\begin{theorem}[Existence and Uniqueness of the Empty Set]\label{thm:empty-set-exists-unique}
There exists a unique set $E$ such that no object is an element of $E$:
\[
  \exists! E\,\forall x,\ x\notin E.
\]
\hyperref[prf:empty-set-exists-unique]{Go to proof.}
\end{theorem}
\end{theorembox}

\LeanFormalizes{thm:empty-set-exists-unique}{lra-lean}{LRA.VolumeI.Set.TT.Set}{LRASet.EmptySetExistsUnique}{pending}


\begin{remark*}[Standard quantified statement]
\[
  \exists E\,
  \Bigl(\forall x,\ x\notin E\ \wedge\
  \forall E'\,((\forall x,\ x\notin E')\to E'=E)\Bigr).
\]
\end{remark*}

\begin{remark*}[Predicate reading]
\[
  \text{unique realization of }\quad \forall x,\ x\notin E.
\]
\end{remark*}

\begin{remark*}[Interpretation]
The symbol \(\varnothing\) is licensed only after one proves that the
empty-set condition determines exactly one set.
\end{remark*}

\begin{dependencies}
\begin{itemize}
  \item \hyperref[ax:empty-set]{Axiom of Empty Set}
  \item \hyperref[ax:extensionality]{Axiom of Extensionality}
\end{itemize}
\end{dependencies}

% Empty Set
% ---------------------------------------------------------
\begin{definitionbox}{Definition (Empty Set)}
\begin{definition}[Empty Set]\label{def:empty-set}
The unique set with no elements is called the \emph{empty set} and is denoted
$\varnothing$.
\end{definition}
\end{definitionbox}

\begin{remark*}[Standard quantified statement]
\[
  \operatorname{EmptySet}(\varnothing)
  \iff
  \forall x,\ x\notin\varnothing.
\]
\end{remark*}

\begin{remark*}[Predicate reading]
\[
  \operatorname{EmptySet}(\varnothing)
  \iff
  \forall x,\ x\notin\varnothing.
\]
\end{remark*}

\begin{remark*}[Interpretation]
The empty set is the set with no members. It is unique by extensionality.
\end{remark*}

\begin{remark*}[Exposition]
Statements of the form $\forall x \in \varnothing,\; P(x)$ are
\emph{vacuously true}: there is no element $x \in \varnothing$ for which
$P(x)$ could fail. This is not a special convention but a consequence of
how universal quantification is defined.
\end{remark*}

\begin{dependencies}
\begin{itemize}
  \item \hyperref[thm:empty-set-exists-unique]{Existence and Uniqueness of the Empty Set}
\end{itemize}
\end{dependencies}

% ---------------------------------------------------------
% Subset and Set Equality
% ---------------------------------------------------------
\begin{definitionbox}{Definition (Subset)}
\begin{definition}[Subset]\label{def:subset}
Let $A$ and $B$ be sets. We say $A$ is a \emph{subset} of $B$, written
$A \subseteq B$, if every element of $A$ is also an element of $B$:
\[
A \subseteq B \;\;\Longleftrightarrow\;\; \forall x \, (x \in A \rightarrow x \in B).
\]
\end{definition}
\end{definitionbox}

\begin{remark*}[Standard quantified statement]
\[
  A\subseteq B
  \iff
  \forall x,\ (x\in A\to x\in B).
\]
\end{remark*}

\begin{remark*}[Predicate reading]
\[
  \operatorname{Subset}(A,B)
  \iff
  \forall x,\ (x\in A\to x\in B).
\]
\end{remark*}

\begin{remark*}[Interpretation]
Subset inclusion means that membership in $A$ is a sufficient condition for
membership in $B$.
\end{remark*}

\begin{dependencies}
\begin{itemize}
  \item \hyperref[def:universal-q]{Universal Formula}
\end{itemize}
\end{dependencies}

\begin{definitionbox}{Definition (Proper Subset)}
\begin{definition}[Proper Subset]\label{def:proper-subset}
$A$ is a \emph{proper subset} of $B$, written $A \subsetneq B$, if
$A \subseteq B$ and $A \neq B$.
\end{definition}
\end{definitionbox}

\begin{remark*}[Standard quantified statement]
\[
  A\subsetneq B
  \iff
  A\subseteq B \wedge A\neq B.
\]
\end{remark*}

\begin{remark*}[Predicate reading]
\[
  \operatorname{ProperSubset}(A,B)
  \iff
  \operatorname{Subset}(A,B)\wedge \neg\operatorname{SetEqual}(A,B).
\]
\end{remark*}

\begin{remark*}[Interpretation]
A proper subset is included in the larger set but does not exhaust it.
\end{remark*}

\begin{remark*}[Exposition]
Some authors write $A \subset B$ for a proper subset; others use it to mean
$A \subseteq B$. In these notes, $\subseteq$ always denotes subset (possibly
equal) and $\subsetneq$ always denotes proper subset.
\end{remark*}

\begin{dependencies}
\begin{itemize}
  \item \hyperref[def:subset]{Subset}
\end{itemize}
\end{dependencies}

\begin{definitionbox}{Definition (Set Equality)}
\begin{definition}[Set Equality]\label{def:set-equality}
Two sets $A$ and $B$ are \emph{equal}, written $A = B$, if they have the same
elements:
\[
A = B \;\;\Longleftrightarrow\;\; \forall x \, (x \in A \leftrightarrow x \in B).
\]
Equivalently, $A = B \iff (A \subseteq B \land B \subseteq A)$.
\end{definition}
\end{definitionbox}

\begin{remark*}[Standard quantified statement]
\[
  A=B
  \iff
  \forall x,\ (x\in A\leftrightarrow x\in B).
\]
\end{remark*}

\begin{remark*}[Predicate reading]
\[
  \operatorname{SetEqual}(A,B)
  \iff
  \operatorname{Subset}(A,B)\wedge \operatorname{Subset}(B,A).
\]
\end{remark*}

\begin{remark*}[Interpretation]
Set equality is extensional: two sets are equal exactly when they have the
same elements.
\end{remark*}

\begin{remark*}[Exposition]
In practice, to prove $A = B$ one shows mutual inclusion: first $A \subseteq B$
(take arbitrary $x \in A$ and deduce $x \in B$), then $B \subseteq A$.
This two-step structure appears throughout set-theoretic proofs.
\end{remark*}

\begin{dependencies}
\begin{itemize}
  \item \hyperref[def:subset]{Subset}
  \item \hyperref[ax:extensionality]{Axiom of Extensionality}
\end{itemize}
\end{dependencies}

\begin{theorembox}{Theorem (Extensional Equality Criteria)}
\begin{theorem}[Extensional Equality Criteria]\label{thm:extensional-equality-criteria}
Let $A$ and $B$ be sets. Then
\[
A=B \iff \forall x\,(x\in A\leftrightarrow x\in B),
\]
and
\[
A=B \iff (A\subseteq B \land B\subseteq A).
\]
\hyperref[prf:extensional-equality-criteria]{Go to proof.}
\end{theorem}
\end{theorembox}

\begin{remark*}[Standard quantified statement]
\[
  A=B \iff \forall x\,(x\in A\leftrightarrow x\in B),
  \qquad
  A=B \iff (A\subseteq B \wedge B\subseteq A).
\]
\end{remark*}

\begin{remark*}[Predicate reading]
\[
  \operatorname{SetEqual}(A,B)
  \iff
  \operatorname{Subset}(A,B)\wedge\operatorname{Subset}(B,A).
\]
\end{remark*}

\begin{remark*}[Interpretation]
To prove equality of sets, prove that they have the same members. In most
working proofs this becomes the mutual-inclusion method: prove $A\subseteq B$
and $B\subseteq A$.
\end{remark*}

\LeanFormalizes{thm:extensional-equality-criteria}{lra-lean}{LRA.VolumeI.Set.Operations.Laws.SubsetCriteria}{SetExtensionalityIff}{pending}
\LeanFormalizes{thm:extensional-equality-criteria}{lra-lean}{LRA.VolumeI.Set.Operations.Laws.SubsetCriteria}{SetEqualityIffMutualSubset}{pending}

\begin{dependencies}
\begin{itemize}
  \item \hyperref[def:subset]{Subset}
  \item \hyperref[def:set-equality]{Set Equality}
  \item \hyperref[ax:extensionality]{Axiom of Extensionality}
\end{itemize}
\end{dependencies}

\begin{theorembox}{Theorem (Subset Absorption Criteria)}
\begin{theorem}[Subset Absorption Criteria]\label{thm:subset-absorption-criteria}
Let $A$ and $B$ be sets. Then
\[
A\subseteq B \iff A\cup B=B,
\]
and
\[
A\subseteq B \iff A\cap B=A.
\]
\hyperref[prf:subset-absorption-criteria]{Go to proof.}
\end{theorem}
\end{theorembox}

\begin{remark*}[Standard quantified statement]
\[
  \forall A,B,\quad
  A\subseteq B\iff A\cup B=B,
  \qquad
  A\subseteq B\iff A\cap B=A.
\]
\end{remark*}

\begin{remark*}[Predicate reading]
\[
  \operatorname{Subset}(A,B)
  \iff
  \operatorname{SetEqual}(A\cup B,B),
  \qquad
  \operatorname{Subset}(A,B)
  \iff
  \operatorname{SetEqual}(A\cap B,A).
\]
\end{remark*}

\begin{remark*}[Interpretation]
Inclusion can be recognized algebraically: a smaller set is absorbed by union
with the larger set and by intersection with the larger set.
\end{remark*}

\LeanFormalizes{thm:subset-absorption-criteria}{lra-lean}{LRA.VolumeI.Set.Operations.Laws.SubsetCriteria}{SubsetIffUnionEqRight}{pending}
\LeanFormalizes{thm:subset-absorption-criteria}{lra-lean}{LRA.VolumeI.Set.Operations.Laws.SubsetCriteria}{SubsetIffIntersectionEqLeft}{pending}

\begin{dependencies}
\begin{itemize}
  \item \hyperref[def:subset]{Subset}
  \item \hyperref[def:set-equality]{Set Equality}
  \item \hyperref[def:union]{Union}
  \item \hyperref[def:intersection]{Intersection}
  \item \hyperref[ax:extensionality]{Axiom of Extensionality}
\end{itemize}
\end{dependencies}

% ---------------------------------------------------------
\begin{theorembox}{Theorem (Existence and Uniqueness of Pairing's Output)}
\begin{theorem}[Existence and Uniqueness of Pairing's Output]\label{thm:pairing-output-exists-unique}
For any sets $x_1$ and $x_2$, there exists a unique set $P$ whose elements are
exactly $x_1$ and $x_2$:
\[
  \forall x_1\,\forall x_2\,\exists! P\,\forall w,\quad
  w\in P\leftrightarrow (w=x_1\vee w=x_2).
\]
\hyperref[prf:pairing-output-exists-unique]{Go to proof.}
\end{theorem}
\end{theorembox}

\LeanFormalizes{thm:pairing-output-exists-unique}{lra-lean}{LRA.VolumeI.Set.TT.Set}{LRASet.PairSetExistsUnique}{pending}


\begin{remark*}[Standard quantified statement]
\[
  \forall x_1\,\forall x_2\,\exists P\,
  \Bigl(
  \forall w,\ (w\in P\leftrightarrow(w=x_1\vee w=x_2))
  \wedge
  \forall P'\,
  ((\forall w,\ (w\in P'\leftrightarrow(w=x_1\vee w=x_2)))\to P'=P)
  \Bigr).
\]
\end{remark*}

\begin{remark*}[Predicate reading]
\[
  \text{unique realization of }\quad
  \forall w,\ (w\in P\leftrightarrow(w=x_1\vee w=x_2)).
\]
\end{remark*}

\begin{remark*}[Interpretation]
For fixed inputs \(x_1\) and \(x_2\), the brace notation
\(\{x_1,x_2\}\) names a unique output set. The singleton \(\{x\}\) is the
special case \(\{x,x\}\), whose only member is \(x\).
\end{remark*}

\begin{dependencies}
\begin{itemize}
  \item \hyperref[ax:pairing]{Axiom of Pairing}
  \item \hyperref[ax:extensionality]{Axiom of Extensionality}
\end{itemize}
\end{dependencies}

\begin{theorembox}{Theorem (Existence and Uniqueness of Union's Output)}
\begin{theorem}[Existence and Uniqueness of Union's Output]\label{thm:union-output-exists-unique}
For any set $x$, there exists a unique set $U$ whose elements are exactly the
elements of members of $x$:
\[
  \forall x\,\exists! U\,\forall z,\quad
  z\in U\leftrightarrow \exists w\,(w\in x\wedge z\in w).
\]
\hyperref[prf:union-output-exists-unique]{Go to proof.}
\end{theorem}
\end{theorembox}

\LeanFormalizes{thm:union-output-exists-unique}{lra-lean}{LRA.VolumeI.Set.Families}{IndexedUnionExistsUnique}{pending}


\begin{remark*}[Standard quantified statement]
\[
  \forall x\,\exists U\,
  \Bigl(
  \forall z,\ (z\in U\leftrightarrow \exists w\,(w\in x\wedge z\in w))
  \wedge
  \forall U'\,
  ((\forall z,\ (z\in U'\leftrightarrow \exists w\,(w\in x\wedge z\in w)))
  \to U'=U)
  \Bigr).
\]
\end{remark*}

\begin{remark*}[Predicate reading]
\[
  \text{unique realization of }\quad
  \forall z,\ (z\in U\leftrightarrow\exists w\,(w\in x\wedge z\in w)).
\]
\end{remark*}

\begin{remark*}[Interpretation]
For a set-indexed family \(x\), the notation \(\bigcup x\) names the unique
set collecting elements of members of \(x\).
\end{remark*}

\begin{dependencies}
\begin{itemize}
  \item \hyperref[ax:union]{Axiom of Union}
  \item \hyperref[ax:extensionality]{Axiom of Extensionality}
\end{itemize}
\end{dependencies}

\begin{corollarybox}{Corollary (Existence and Uniqueness of Binary Union)}
\begin{corollary}[Existence and Uniqueness of Binary Union]\label{cor:binary-union-exists-unique}
For any sets $A$ and $B$, the set
\[
  A\cup B:=\bigcup\{A,B\}
\]
exists and is unique.
\hyperref[prf:binary-union-exists-unique]{Go to proof.}
\end{corollary}
\end{corollarybox}

\LeanFormalizes{cor:binary-union-exists-unique}{lra-lean}{LRA.VolumeI.Set.TT.Set}{LRASet.BinaryUnionExistsUnique}{pending}


\begin{remark*}[Standard quantified statement]
\[
  \forall A\,\forall B\,\exists! C\,\forall x,\quad
  x\in C\leftrightarrow(x\in A\vee x\in B).
\]
\end{remark*}

\begin{remark*}[Predicate reading]
\[
  \text{unique realization of }\quad
  \forall x,\ (x\in C\leftrightarrow(x\in A\vee x\in B)).
\]
\end{remark*}

\begin{remark*}[Interpretation]
Binary union is a derived operation obtained by first pairing \(A,B\) and then
taking the union of that pair set.
\end{remark*}

\begin{dependencies}
\begin{itemize}
  \item \hyperref[thm:pairing-output-exists-unique]{Existence and Uniqueness of Pairing's Output}
  \item \hyperref[thm:union-output-exists-unique]{Existence and Uniqueness of Union's Output}
\end{itemize}
\end{dependencies}

% Set Operations
% ---------------------------------------------------------
\begin{lratopiccard}{Basic Operations --- Membership Signatures}
\begin{lrasignaturetable}
\LraSignatureRow{Context}{Work inside a fixed ambient set \(U\) when complements are used.}
\LraSignatureRow{Binary union}{\(A\cup B\): \(x\in A\cup B\) iff \(x\in A\) or \(x\in B\).}
\LraSignatureRow{Binary intersection}{\(A\cap B\): \(x\in A\cap B\) iff \(x\in A\) and \(x\in B\).}
\LraSignatureRow{Difference}{\(A\setminus B\): \(x\in A\setminus B\) iff \(x\in A\) and \(x\notin B\).}
\LraSignatureRow{Symmetric difference}{\(A\triangle B\): membership means \(x\) belongs to exactly one of \(A,B\).}
\LraSignatureRow{Complement}{\(A^c\): \(x\in A^c\) iff \(x\in U\) and \(x\notin A\).}
\LraSignatureRow{Power set}{\(\mathcal P(A)\): \(S\in\mathcal P(A)\) iff \(S\subseteq A\).}
\end{lrasignaturetable}
\end{lratopiccard}

\begin{definitionbox}{Definition (Union)}
\begin{definition}[Union]\label{def:union}
The \emph{union} of $A$ and $B$, denoted $A \cup B$, is the set of all
elements belonging to at least one of the two sets:
\[
A \cup B \;=\; \{\, x \mid x \in A \lor x \in B \,\}.
\]
\end{definition}
\end{definitionbox}

\begin{remark*}[Standard quantified statement]
\[
  \forall x,\quad x\in A\cup B \iff (x\in A\vee x\in B).
\]
\end{remark*}

\begin{remark*}[Predicate reading]
\[
  \operatorname{UnionSet}(C,A,B)
  \iff
  \forall x,\ (x\in C \leftrightarrow (x\in A\vee x\in B)).
\]
\end{remark*}

\begin{remark*}[Interpretation]
The union keeps every element that belongs to at least one of the two sets.
\end{remark*}

\begin{dependencies}
\begin{itemize}
  \item \hyperref[def:set-membership]{Set and Membership}
  \item \hyperref[cor:binary-union-exists-unique]{Existence and Uniqueness of Binary Union}
\end{itemize}
\end{dependencies}

\begin{lracategoricalcard}{Categorical Rendering --- Union as Join}
\begin{center}
\begin{tikzpicture}[lra categorical diagram]
  \node[lra cat native object] (A) at (0,1.65)
    {\LraCatObject{A}{\{x:x\in A\}}};
  \node[lra cat native object] (B) at (4.2,1.65)
    {\LraCatObject{B}{\{x:x\in B\}}};

  \node[lra cat result object] (AuB) at (2.1,0)
    {\LraCatObject{A\cup B}{\{x:x\in A\lor x\in B\}}};

  \node[lra cat native object] (C) at (2.1,-1.65)
    {\LraCatObject{C}{A\subseteq C,\ B\subseteq C}};

  \draw[lra cat left arrow] (A) -- node[lra cat label,above left]{include} (AuB);
  \draw[lra cat right arrow] (B) -- node[lra cat label,above right]{include} (AuB);
  \draw[lra cat universal arrow] (AuB) -- node[lra cat label,right]{least upper bound} (C);

  \node[lra cat note] at (2.1,-2.35)
    {\(A\cup B\) is the smallest set containing both \(A\) and \(B\)};
\end{tikzpicture}

\end{center}
\[
  A\subseteq C\ \wedge\ B\subseteq C
  \quad\Longleftrightarrow\quad
  A\cup B\subseteq C.
\]
\end{lracategoricalcard}

\begin{lracategoricalcard}{Concrete Picture --- Union}
\begin{center}
\begin{tikzpicture}[atlas, scale=0.85]
  \fill[atlasblue!18] (-1,0) circle (1.5);
  \fill[atlasblue!18] ( 1,0) circle (1.5);
  \draw[atlasblue, thick] (-1,0) circle (1.5);
  \draw[atlasblue, thick] ( 1,0) circle (1.5);
  \node[atlas label] at (-1,0.05) {$A$};
  \node[atlas label] at ( 1,0.05) {$B$};
  \node[atlas label] at (0,-1.95) {$A\cup B=\{x:x\in A\lor x\in B\}$};
\end{tikzpicture}

\end{center}
\[
  x\in A\cup B
  \quad\Longleftrightarrow\quad
  x\in A\ \vee\ x\in B.
\]
\end{lracategoricalcard}

\begin{theorembox}{Theorem (Existence and Uniqueness of Separation's Output)}
\begin{theorem}[Existence and Uniqueness of Separation's Output]\label{thm:separation-output-exists-unique}
Let $A$ be a set and let $\varphi(x)$ be a formula. There exists a unique set
$B$ whose elements are exactly the elements of $A$ satisfying $\varphi$:
\[
  \forall A\,\exists! B\,\forall x,\quad
  x\in B\leftrightarrow (x\in A\wedge \varphi(x)).
\]
\hyperref[prf:separation-output-exists-unique]{Go to proof.}
\end{theorem}
\end{theorembox}

\LeanFormalizes{thm:separation-output-exists-unique}{lra-lean}{LRA.VolumeI.Set.TT.Set}{LRASet.SeparationExistsUnique}{pending}


\begin{remark*}[Standard quantified statement]
\[
  \forall A\,\exists B\,
  \Bigl(
  \forall x,\ (x\in B\leftrightarrow(x\in A\wedge\varphi(x)))
  \wedge
  \forall B'\,
  ((\forall x,\ (x\in B'\leftrightarrow(x\in A\wedge\varphi(x))))\to B'=B)
  \Bigr).
\]
\end{remark*}

\begin{remark*}[Predicate reading]
\[
  \text{unique realization of }\quad
  \forall x,\ (x\in B\leftrightarrow(x\in A\wedge\varphi(x))).
\]
\end{remark*}

\begin{remark*}[Interpretation]
Bounded set-builder notation \(\{x\in A\mid \varphi(x)\}\) is legitimate
because Separation supplies a witness and Extensionality makes it unique.
\end{remark*}

\begin{dependencies}
\begin{itemize}
  \item \hyperref[ax:separation]{Axiom Schema of Separation}
  \item \hyperref[ax:extensionality]{Axiom of Extensionality}
\end{itemize}
\end{dependencies}

\begin{corollarybox}{Corollary (Existence and Uniqueness of Intersection)}
\begin{corollary}[Existence and Uniqueness of Intersection]\label{cor:intersection-exists-unique}
For any sets $A$ and $B$, the set
\[
  A\cap B:=\{\,x\in A\mid x\in B\,\}
\]
exists and is unique.
\hyperref[prf:intersection-exists-unique]{Go to proof.}
\end{corollary}
\end{corollarybox}

\LeanFormalizes{cor:intersection-exists-unique}{lra-lean}{LRA.VolumeI.Set.TT.Set}{LRASet.IntersectionExistsUnique}{pending}


\begin{remark*}[Standard quantified statement]
\[
  \forall A\,\forall B\,\exists! C\,\forall x,\quad
  x\in C\leftrightarrow(x\in A\wedge x\in B).
\]
\end{remark*}

\begin{remark*}[Predicate reading]
\[
  \text{unique realization of }\quad
  \forall x,\ (x\in C\leftrightarrow(x\in A\wedge x\in B)).
\]
\end{remark*}

\begin{remark*}[Interpretation]
Intersection is the separated subset of \(A\) consisting of elements that also
belong to \(B\).
\end{remark*}

\begin{dependencies}
\begin{itemize}
  \item \hyperref[thm:separation-output-exists-unique]{Existence and Uniqueness of Separation's Output}
\end{itemize}
\end{dependencies}

\begin{definitionbox}{Definition (Intersection)}
\begin{definition}[Intersection]\label{def:intersection}
The \emph{intersection} of $A$ and $B$, denoted $A \cap B$, is the set of all
elements common to both:
\[
A \cap B \;=\; \{\, x \mid x \in A \land x \in B \,\}.
\]
\end{definition}
\end{definitionbox}

\begin{remark*}[Standard quantified statement]
\[
  \forall x,\quad x\in A\cap B \iff (x\in A\wedge x\in B).
\]
\end{remark*}

\begin{remark*}[Predicate reading]
\[
  \operatorname{IntersectionSet}(C,A,B)
  \iff
  \forall x,\ (x\in C \leftrightarrow (x\in A\wedge x\in B)).
\]
\end{remark*}

\begin{remark*}[Interpretation]
The intersection keeps exactly the elements common to both sets.
\end{remark*}

\begin{dependencies}
\begin{itemize}
  \item \hyperref[def:set-membership]{Set and Membership}
  \item \hyperref[cor:intersection-exists-unique]{Existence and Uniqueness of Intersection}
\end{itemize}
\end{dependencies}

\begin{lracategoricalcard}{Categorical Rendering --- Intersection as Meet}
\begin{center}
\begin{tikzpicture}[lra categorical diagram]
  \node[lra cat native object] (A) at (0,1.65)
    {\LraCatObject{A}{\{x:x\in A\}}};
  \node[lra cat native object] (B) at (4.2,1.65)
    {\LraCatObject{B}{\{x:x\in B\}}};

  \node[lra cat result object] (AnB) at (2.1,0)
    {\LraCatObject{A\cap B}{\{x:x\in A\land x\in B\}}};

  \node[lra cat native object] (D) at (2.1,-1.65)
    {\LraCatObject{D}{D\subseteq A,\ D\subseteq B}};

  \draw[lra cat result arrow] (AnB) -- node[lra cat label,above left]{include} (A);
  \draw[lra cat result arrow] (AnB) -- node[lra cat label,above right]{include} (B);
  \draw[lra cat universal arrow] (D) -- node[lra cat label,right]{greatest lower bound} (AnB);

  \node[lra cat note] at (2.1,-2.35)
    {\(A\cap B\) is the largest set contained in both \(A\) and \(B\)};
\end{tikzpicture}

\end{center}
\[
  D\subseteq A\ \wedge\ D\subseteq B
  \quad\Longleftrightarrow\quad
  D\subseteq A\cap B.
\]
\end{lracategoricalcard}

\begin{lracategoricalcard}{Concrete Picture --- Intersection}
\begin{center}
\begin{tikzpicture}[atlas, scale=0.85]
  \begin{scope}
    \clip (-1,0) circle (1.5);
    \fill[atlasgreen!24] (1,0) circle (1.5);
  \end{scope}
  \draw[atlasblue, thick] (-1,0) circle (1.5);
  \draw[atlasblue, thick] ( 1,0) circle (1.5);
  \node[atlas label] at (-1.35,0.05) {$A$};
  \node[atlas label] at ( 1.35,0.05) {$B$};
  \node[atlas label] at (0,-1.95) {$A\cap B=\{x:x\in A\land x\in B\}$};
\end{tikzpicture}

\end{center}
\[
  x\in A\cap B
  \quad\Longleftrightarrow\quad
  x\in A\ \wedge\ x\in B.
\]
\end{lracategoricalcard}

\begin{corollarybox}{Corollary (Existence and Uniqueness of Set Difference)}
\begin{corollary}[Existence and Uniqueness of Set Difference]\label{cor:set-difference-exists-unique}
For any sets $A$ and $B$, the set
\[
  A\setminus B:=\{\,x\in A\mid x\notin B\,\}
\]
exists and is unique.
\hyperref[prf:set-difference-exists-unique]{Go to proof.}
\end{corollary}
\end{corollarybox}

\LeanFormalizes{cor:set-difference-exists-unique}{lra-lean}{LRA.VolumeI.Set.TT.Set}{LRASet.DifferenceExistsUnique}{pending}


\begin{remark*}[Standard quantified statement]
\[
  \forall A\,\forall B\,\exists! C\,\forall x,\quad
  x\in C\leftrightarrow(x\in A\wedge x\notin B).
\]
\end{remark*}

\begin{remark*}[Predicate reading]
\[
  \text{unique realization of }\quad
  \forall x,\ (x\in C\leftrightarrow(x\in A\wedge x\notin B)).
\]
\end{remark*}

\begin{remark*}[Interpretation]
Set difference is the separated subset of \(A\) consisting of elements outside
\(B\).
\end{remark*}

\begin{dependencies}
\begin{itemize}
  \item \hyperref[thm:separation-output-exists-unique]{Existence and Uniqueness of Separation's Output}
\end{itemize}
\end{dependencies}

\begin{definitionbox}{Definition (Set Difference)}
\begin{definition}[Set Difference]\label{def:set-difference}
The \emph{set difference} $A \setminus B$ is the set of elements in $A$ but
not in $B$:
\[
A \setminus B \;=\; \{\, x \mid x \in A \land x \notin B \,\}.
\]
\end{definition}
\end{definitionbox}

\begin{remark*}[Standard quantified statement]
\[
  \forall x,\quad x\in A\setminus B \iff (x\in A\wedge x\notin B).
\]
\end{remark*}

\begin{remark*}[Predicate reading]
\[
  \operatorname{SetDifference}(C,A,B)
  \iff
  \forall x,\ (x\in C \leftrightarrow (x\in A\wedge x\notin B)).
\]
\end{remark*}

\begin{remark*}[Interpretation]
Set difference removes from $A$ every element that also belongs to $B$.
\end{remark*}

\begin{dependencies}
\begin{itemize}
  \item \hyperref[def:set-membership]{Set and Membership}
  \item \hyperref[cor:set-difference-exists-unique]{Existence and Uniqueness of Set Difference}
\end{itemize}
\end{dependencies}

\begin{center}
\begin{tikzpicture}[scale=0.85]
  \fill[pattern=north east lines] (-1,0) circle (1.5);
  \begin{scope}
    \clip (1,0) circle (1.5);
    \fill[white] (-1,0) circle (1.5);
  \end{scope}
  \draw (-1,0) circle (1.5);
  \draw ( 1,0) circle (1.5);
  \node at (-2.2,-1.6) {$A$};
  \node at ( 2.2,-1.6) {$B$};
\end{tikzpicture}
\end{center}


\begin{corollarybox}{Corollary (Existence and Uniqueness of Symmetric Difference)}
\begin{corollary}[Existence and Uniqueness of Symmetric Difference]\label{cor:symmetric-difference-exists-unique}
For any sets $A$ and $B$, the set
\[
  A\triangle B:=(A\setminus B)\cup(B\setminus A)
\]
exists and is unique.
\hyperref[prf:symmetric-difference-exists-unique]{Go to proof.}
\end{corollary}
\end{corollarybox}

\LeanFormalizes{cor:symmetric-difference-exists-unique}{lra-lean}{LRA.VolumeI.Set.TT.Set}{LRASet.SymmetricDifferenceExistsUnique}{pending}


\begin{remark*}[Standard quantified statement]
\[
  \forall A\,\forall B\,\exists! C\,\forall x,\quad
  x\in C\leftrightarrow
  ((x\in A\wedge x\notin B)\vee(x\in B\wedge x\notin A)).
\]
\end{remark*}

\begin{remark*}[Predicate reading]
\[
  \text{unique realization of }\quad
  \forall x,\ (x\in C\leftrightarrow((x\in A\wedge x\notin B)\vee
  (x\in B\wedge x\notin A))).
\]
\end{remark*}

\begin{remark*}[Interpretation]
Symmetric difference is a composite set-former built from two set differences
and one binary union.
\end{remark*}

\begin{dependencies}
\begin{itemize}
  \item \hyperref[cor:set-difference-exists-unique]{Existence and Uniqueness of Set Difference}
  \item \hyperref[cor:binary-union-exists-unique]{Existence and Uniqueness of Binary Union}
\end{itemize}
\end{dependencies}

\begin{definitionbox}{Definition (Symmetric Difference)}
\begin{definition}[Symmetric Difference]\label{def:sym-diff}
The \emph{symmetric difference} $A \triangle B$ is the set of elements
belonging to exactly one of $A$ and $B$:
\[
A \triangle B
\;=\;
(A \setminus B) \cup (B \setminus A)
\;=\;
\{\, x \mid (x \in A \lor x \in B) \land x \notin A \cap B \,\}.
\]
\end{definition}
\end{definitionbox}

\begin{remark*}[Standard quantified statement]
\[
  \forall x,\quad x\in A\triangle B
  \iff
  ((x\in A\wedge x\notin B)\vee(x\in B\wedge x\notin A)).
\]
\end{remark*}

\begin{remark*}[Predicate reading]
\[
  \operatorname{SymmetricDifference}(C,A,B)
  \iff
  \forall x,\ (x\in C \leftrightarrow ((x\in A\wedge x\notin B)\vee(x\in B\wedge x\notin A))).
\]
\end{remark*}

\begin{remark*}[Interpretation]
The symmetric difference keeps elements that belong to exactly one of the two
sets.
\end{remark*}

\begin{dependencies}
\begin{itemize}
  \item \hyperref[cor:symmetric-difference-exists-unique]{Existence and Uniqueness of Symmetric Difference}
\end{itemize}
\end{dependencies}

\input{volume-i/book-sets/set-theory/notes/sets/figure-symmetric-difference}

\begin{corollarybox}{Corollary (Existence and Uniqueness of Relative Complement)}
\begin{corollary}[Existence and Uniqueness of Relative Complement]\label{cor:relative-complement-exists-unique}
Let $U$ be a fixed ambient set. For every set $A$, the relative complement
\[
  A^c:=U\setminus A
\]
exists and is unique. This is a relative construction; ZFC has no absolute
complement operation over a universal set of all sets.
\hyperref[prf:relative-complement-exists-unique]{Go to proof.}
\end{corollary}
\end{corollarybox}

\LeanFormalizes{cor:relative-complement-exists-unique}{lra-lean}{LRA.VolumeI.Set.TT.Set}{LRASet.RelativeComplementExistsUnique}{pending}


\begin{remark*}[Standard quantified statement]
\[
  \forall A\,\exists! C\,\forall x,\quad
  x\in C\leftrightarrow(x\in U\wedge x\notin A).
\]
\end{remark*}

\begin{remark*}[Predicate reading]
\[
  \text{unique realization of }\quad
  \forall x,\ (x\in C\leftrightarrow(x\in U\wedge x\notin A)).
\]
\end{remark*}

\begin{remark*}[Interpretation]
Complement notation is relative to a fixed ambient set \(U\), not an absolute
operation over all sets.
\end{remark*}

\begin{dependencies}
\begin{itemize}
  \item \hyperref[cor:set-difference-exists-unique]{Existence and Uniqueness of Set Difference}
\end{itemize}
\end{dependencies}

\begin{definitionbox}{Definition (Complement)}
\begin{definition}[Complement]\label{def:complement}
Let $U$ be a fixed universe and $A \subseteq U$. The \emph{complement} of $A$
relative to $U$, denoted $A^c$, is:
\[
A^c \;=\; U \setminus A.
\]
\end{definition}
\end{definitionbox}

\begin{remark*}[Standard quantified statement]
\[
  \forall x,\quad x\in A^c \iff (x\in U\wedge x\notin A).
\]
\end{remark*}

\begin{remark*}[Predicate reading]
\[
  \operatorname{ComplementSet}(C,A,U)
  \iff
  \forall x,\ (x\in C \leftrightarrow (x\in U\wedge x\notin A)).
\]
\end{remark*}

\begin{remark*}[Interpretation]
The complement of $A$ depends on the ambient universe $U$; it contains the
elements of $U$ that are outside $A$.
\end{remark*}

\begin{dependencies}
\begin{itemize}
  \item \hyperref[def:set-membership]{Set and Membership}
  \item \hyperref[cor:relative-complement-exists-unique]{Existence and Uniqueness of Relative Complement}
\end{itemize}
\end{dependencies}

\begin{definitionbox}{Definition (Cartesian Product)}
\begin{definition}[Cartesian Product]\label{def:cartesian-product}
The \emph{Cartesian product} of sets $A$ and $B$ is:
\[
A \times B = \{\, (a,b) \mid a \in A \text{ and } b \in B \,\}.
\]
\end{definition}
\end{definitionbox}

\begin{remark*}[Standard quantified statement]
\[
  \forall x,\quad x\in A\times B
  \iff
  \exists a\in A\,\exists b\in B,\ x=(a,b).
\]
\end{remark*}

\begin{remark*}[Predicate reading]
\[
  \operatorname{CartesianProduct}(P,A,B)
  \iff
  \forall x,\ (x\in P\leftrightarrow \exists a\in A\,\exists b\in B,\ x=(a,b)).
\]
\end{remark*}

\begin{remark*}[Interpretation]
The Cartesian product records ordered choices of one element from $A$ and one
element from $B$.
\end{remark*}

\begin{remark*}[Exposition]
The ordered pair $(a,b)$ is formally defined as $\{\{a\},\{a,b\}\}$.
This Kuratowski encoding ensures the key property: $(a,b) = (c,d)$ iff
$a = c$ and $b = d$.
\end{remark*}

\begin{remark*}[Exposition]
The Cartesian product provides a way to encode ordered information. Relations
and functions will be defined as special subsets of Cartesian products, so they
require no new foundational objects.
\end{remark*}

\begin{dependencies}
\begin{itemize}
  \item \hyperref[def:ordered-pair]{Ordered Pair}
  \item \hyperref[ax:pairing]{Axiom of Pairing}
  \item \hyperref[ax:power-set]{Axiom of Power Set}
  \item \hyperref[ax:separation]{Axiom Schema of Separation}
\end{itemize}
\end{dependencies}

\begin{theorembox}{Theorem (Existence and Uniqueness of Power Set's Output)}
\begin{theorem}[Existence and Uniqueness of Power Set's Output]\label{thm:power-set-output-exists-unique}
For any set $A$, there exists a unique set $P$ whose elements are exactly the
subsets of $A$:
\[
  \forall A\,\exists! P\,\forall S,\quad
  S\in P\leftrightarrow S\subseteq A.
\]
\hyperref[prf:power-set-output-exists-unique]{Go to proof.}
\end{theorem}
\end{theorembox}

\LeanFormalizes{thm:power-set-output-exists-unique}{lra-lean}{LRA.VolumeI.Set.PowerSets}{RelativePowerSetExistsUnique}{pending}


\begin{remark*}[Standard quantified statement]
\[
  \forall A\,\exists P\,
  \Bigl(
  \forall S,\ (S\in P\leftrightarrow S\subseteq A)
  \wedge
  \forall P'\,((\forall S,\ (S\in P'\leftrightarrow S\subseteq A))\to P'=P)
  \Bigr).
\]
\end{remark*}

\begin{remark*}[Predicate reading]
\[
  \text{unique realization of }\quad
  \forall S,\ (S\in P\leftrightarrow S\subseteq A).
\]
\end{remark*}

\begin{remark*}[Interpretation]
For each set \(A\), the notation \(\mathcal P(A)\) names the unique set whose
members are exactly the subsets of \(A\).
\end{remark*}

\begin{dependencies}
\begin{itemize}
  \item \hyperref[ax:power-set]{Axiom of Power Set}
  \item \hyperref[ax:extensionality]{Axiom of Extensionality}
\end{itemize}
\end{dependencies}

\begin{definitionbox}{Definition (Power Set)}
\begin{definition}[Power Set]\label{def:power-set}
The \emph{power set} of $A$, denoted $\mathcal{P}(A)$, is the set of all
subsets of $A$:
\[
\mathcal{P}(A) \;:=\; \{\, S \mid S \subseteq A \,\}.
\]
\end{definition}
\end{definitionbox}

\begin{remark*}[Standard quantified statement]
\[
  \forall S,\quad S\in\mathcal{P}(A)\iff S\subseteq A.
\]
\end{remark*}

\begin{remark*}[Predicate reading]
\[
  \operatorname{PowerSet}(P,A)
  \iff
  \forall S,\ (S\in P\leftrightarrow \operatorname{Subset}(S,A)).
\]
\end{remark*}

\begin{remark*}[Interpretation]
The power set collects every subset of $A$ into a single set.
\end{remark*}

\begin{remark*}[Exposition]
If $A$ has $n$ elements, then $|\mathcal{P}(A)| = 2^n$. For infinite $A$,
Cantor's theorem shows $|\mathcal{P}(A)| > |A|$ strictly.
\end{remark*}

\begin{dependencies}
\begin{itemize}
  \item \hyperref[thm:power-set-output-exists-unique]{Existence and Uniqueness of Power Set's Output}
\end{itemize}
\end{dependencies}

% =========================================================
% The Inclusion Order on Sets
% =========================================================

\subsection{The Inclusion Order on Sets}

% ---------------------------------------------------------
% Monotonicity under inclusion
% ---------------------------------------------------------
\begin{lratopiccard}{Inclusion Order --- Signature}
\begin{lrasignaturetable}
\LraSignatureRow{Relation}{\(A\subseteq B\) iff every member of \(A\) is a member of \(B\).}
\LraSignatureRow{Strict relation}{\(A\subsetneq B\) iff \(A\subseteq B\) and \(A\neq B\).}
\LraSignatureRow{Equality test}{\(A=B\) iff \(A\subseteq B\) and \(B\subseteq A\).}
\LraSignatureRow{Monotone maps}{\(A\subseteq B\Rightarrow F(A)\subseteq F(B)\): the operation preserves inclusion.}
\LraSignatureRow{Antitone maps}{\(A\subseteq B\Rightarrow F(B)\subseteq F(A)\): the operation reverses inclusion.}
\LraSignatureRow{Structural echo}{Sets ordered by inclusion are the first poset/subobject picture in the text.}
\end{lrasignaturetable}
\end{lratopiccard}

\begin{lraproofpatterncard}{Inclusion Proof Patterns}
\begin{lraproofpatterntable}
\LraProofPatternRow{Prove \(A\subseteq B\)}{Take arbitrary \(x\in A\), then prove \(x\in B\).}
\LraProofPatternRow{Prove \(A=B\)}{Prove \(A\subseteq B\) and \(B\subseteq A\), then use Extensionality.}
\LraProofPatternRow{Prove monotonicity}{Assume \(A\subseteq B\), unpack membership in \(F(A)\), use the hypothesis, and rebuild membership in \(F(B)\).}
\LraProofPatternRow{Prove antitonicity}{Assume \(A\subseteq B\), but aim for \(F(B)\subseteq F(A)\).}
\end{lraproofpatterntable}
\end{lraproofpatterncard}

\begin{lracategoricalcard}{Categorical Rendering --- Monotone and Antitone}
\begin{center}
\begin{tikzpicture}[lra categorical diagram]
  \node[lra cat native object] (A) at (0,1.6)
    {\LraCatObject{A}{\{x:x\in A\}}};
  \node[lra cat native object] (B) at (4.2,1.6)
    {\LraCatObject{B}{\{x:x\in B\}}};
  \node[lra cat left object] (FA) at (0,0)
    {\LraCatObject{F(A)}{\text{output on }A}};
  \node[lra cat left object] (FB) at (4.2,0)
    {\LraCatObject{F(B)}{\text{output on }B}};
  \node[lra cat right object] (GB) at (0,-1.6)
    {\LraCatObject{G(B)}{\text{output on }B}};
  \node[lra cat right object] (GA) at (4.2,-1.6)
    {\LraCatObject{G(A)}{\text{output on }A}};

  \draw[lra cat arrow] (A) -- node[lra cat label,above]{\(A\subseteq B\)} (B);
  \draw[lra cat left arrow] (A) -- node[lra cat label,left]{\(F\)} (FA);
  \draw[lra cat left arrow] (B) -- node[lra cat label,right]{\(F\)} (FB);
  \draw[lra cat left arrow] (FA) -- node[lra cat label,above]{preserves} (FB);
  \draw[lra cat right arrow] (B) -- node[lra cat label,right]{\(G\)} (GA);
  \draw[lra cat right arrow] (A) -- node[lra cat label,left]{\(G\)} (GB);
  \draw[lra cat right arrow] (GB) -- node[lra cat label,below]{reverses} (GA);

  \node[lra cat note] at (2.1,-2.35)
    {monotone maps preserve inclusion; antitone maps reverse it};
\end{tikzpicture}

\end{center}
\[
  A\subseteq B
  \quad\Longrightarrow\quad
  F(A)\subseteq F(B),
  \qquad
  A\subseteq B
  \quad\Longrightarrow\quad
  G(B)\subseteq G(A).
\]
\end{lracategoricalcard}

\begin{definitionbox}{Definition (Inclusion-Monotone Set Operation)}
\begin{definition}[Inclusion-Monotone Set Operation]\label{def:inclusion-monotone-set-operation}
A set operation $F$ is \emph{inclusion-monotone} in an argument if replacing
that argument by a larger set can only enlarge the value of the operation.
For a one-variable operation, this means
\[
A\subseteq B \quad\Longrightarrow\quad F(A)\subseteq F(B).
\]
\end{definition}
\end{definitionbox}

\begin{remark*}[Standard quantified statement]
\[
  \forall A,B,\quad A\subseteq B \to F(A)\subseteq F(B).
\]
\end{remark*}

\begin{remark*}[Predicate reading]
\[
  \operatorname{InclusionMonotoneOperation}(F)
  \iff
  \forall A,B,\ (\operatorname{Subset}(A,B)\to \operatorname{Subset}(F(A),F(B))).
\]
\end{remark*}

\begin{remark*}[Interpretation]
An inclusion-monotone operation preserves subset order.
\end{remark*}

\begin{dependencies}
\begin{itemize}
  \item \hyperref[def:subset]{Subset}
\end{itemize}
\end{dependencies}

\begin{definitionbox}{Definition (Inclusion-Antitone Set Operation)}
\begin{definition}[Inclusion-Antitone Set Operation]\label{def:inclusion-antitone-set-operation}
A set operation $F$ is \emph{inclusion-antitone} in an argument if replacing
that argument by a larger set can only shrink the value of the operation:
\[
A\subseteq B \quad\Longrightarrow\quad F(B)\subseteq F(A).
\]
\end{definition}
\end{definitionbox}

\begin{remark*}[Standard quantified statement]
\[
  \forall A,B,\quad A\subseteq B \to F(B)\subseteq F(A).
\]
\end{remark*}

\begin{remark*}[Predicate reading]
\[
  \operatorname{InclusionAntitoneOperation}(F)
  \iff
  \forall A,B,\ (\operatorname{Subset}(A,B)\to \operatorname{Subset}(F(B),F(A))).
\]
\end{remark*}

\begin{remark*}[Interpretation]
An inclusion-antitone operation reverses subset order.
\end{remark*}

\begin{dependencies}
\begin{itemize}
  \item \hyperref[def:subset]{Subset}
\end{itemize}
\end{dependencies}

\begin{theorembox}{Theorem (Union is Inclusion-Monotone)}
\begin{theorem}[Union is Inclusion-Monotone]\label{thm:union-monotone-inclusion}
If $A\subseteq B$, then
\[
A\cup C\subseteq B\cup C
\quad\text{and}\quad
C\cup A\subseteq C\cup B.
\]
\hyperref[prf:union-monotone-inclusion]{Go to proof.}
\end{theorem}
\end{theorembox}

\LeanFormalizes{thm:union-monotone-inclusion}{lra-lean}{LRA.VolumeI.Set.Operations.Laws.Union}{UnionMonotoneInclusion}{pending}

\begin{remark*}[Standard quantified statement]
\[
  \forall A,B,C,\quad
  A\subseteq B\to (A\cup C\subseteq B\cup C\wedge C\cup A\subseteq C\cup B).
\]
\end{remark*}

\begin{remark*}[Predicate reading]
\[
  \operatorname{Subset}(A,B)
  \to
  \operatorname{Subset}(A\cup C,B\cup C)
  \wedge
  \operatorname{Subset}(C\cup A,C\cup B).
\]
\end{remark*}

\begin{remark*}[Interpretation]
Taking a union with a fixed set preserves inclusion in either argument.
\end{remark*}

\begin{remark*}[Exposition]
The first implication reads
\[
\forall x,\quad x\in A\cup C \Rightarrow x\in B\cup C,
\]
assuming $\forall x\,(x\in A\Rightarrow x\in B)$. To disprove such an
inclusion, find a witness $x\in A\cup C$ with $x\notin B\cup C$.
\end{remark*}

\begin{dependencies}
\begin{itemize}
  \item \hyperref[def:subset]{Subset}
  \item \hyperref[def:union]{Union}
\end{itemize}
\end{dependencies}

\begin{theorembox}{Theorem (Intersection is Inclusion-Monotone)}
\begin{theorem}[Intersection is Inclusion-Monotone]\label{thm:intersection-monotone-inclusion}
If $A\subseteq B$, then
\[
A\cap C\subseteq B\cap C
\quad\text{and}\quad
C\cap A\subseteq C\cap B.
\]
\hyperref[prf:intersection-monotone-inclusion]{Go to proof.}
\end{theorem}
\end{theorembox}

\LeanFormalizes{thm:intersection-monotone-inclusion}{lra-lean}{LRA.VolumeI.Set.Operations.Laws.Intersection}{IntersectionMonotoneInclusion}{pending}

\begin{remark*}[Standard quantified statement]
\[
  \forall A,B,C,\quad
  A\subseteq B\to (A\cap C\subseteq B\cap C\wedge C\cap A\subseteq C\cap B).
\]
\end{remark*}

\begin{remark*}[Predicate reading]
\[
  \operatorname{Subset}(A,B)
  \to
  \operatorname{Subset}(A\cap C,B\cap C)
  \wedge
  \operatorname{Subset}(C\cap A,C\cap B).
\]
\end{remark*}

\begin{remark*}[Interpretation]
Taking an intersection with a fixed set preserves inclusion in either argument.
\end{remark*}

\begin{dependencies}
\begin{itemize}
  \item \hyperref[def:subset]{Subset}
  \item \hyperref[def:intersection]{Intersection}
\end{itemize}
\end{dependencies}

\begin{theorembox}{Theorem (Power Set is Inclusion-Monotone)}
\begin{theorem}[Power Set is Inclusion-Monotone]\label{thm:power-set-monotone-inclusion}
If $A\subseteq B$, then
\[
\mathcal{P}(A)\subseteq\mathcal{P}(B).
\]
\hyperref[prf:power-set-monotone-inclusion]{Go to proof.}
\end{theorem}
\end{theorembox}

\LeanFormalizes{thm:power-set-monotone-inclusion}{lra-lean}{LRA.VolumeI.Set.PowerSets}{RelativePowerSetMonotoneInclusion}{pending}

\begin{remark*}[Standard quantified statement]
\[
  \forall A,B,\quad A\subseteq B\to \mathcal{P}(A)\subseteq\mathcal{P}(B).
\]
\end{remark*}

\begin{remark*}[Predicate reading]
\[
  \operatorname{Subset}(A,B)
  \to
  \operatorname{Subset}(\mathcal{P}(A),\mathcal{P}(B)).
\]
\end{remark*}

\begin{remark*}[Interpretation]
Every subset of a smaller set is also a subset of the larger set.
\end{remark*}

\begin{dependencies}
\begin{itemize}
  \item \hyperref[def:subset]{Subset}
  \item \hyperref[def:power-set]{Power Set}
\end{itemize}
\end{dependencies}

\begin{theorembox}{Theorem (Complement is Inclusion-Antitone)}
\begin{theorem}[Complement is Inclusion-Antitone]\label{thm:complement-antitone-inclusion}
Let complements be taken relative to a fixed universe $U$. If
$A\subseteq B\subseteq U$, then
\[
B^c\subseteq A^c.
\]
\hyperref[prf:complement-antitone-inclusion]{Go to proof.}
\end{theorem}
\end{theorembox}

\LeanFormalizes{thm:complement-antitone-inclusion}{lra-lean}{LRA.VolumeI.Set.Operations.Laws.Complement}{ComplementAntitoneInclusion}{pending}

\begin{remark*}[Standard quantified statement]
\[
  \forall A,B\subseteq U,\quad A\subseteq B\to B^c\subseteq A^c.
\]
\end{remark*}

\begin{remark*}[Predicate reading]
\[
  \operatorname{Subset}(A,B)
  \to
  \operatorname{Subset}(B^c,A^c).
\]
\end{remark*}

\begin{remark*}[Interpretation]
When a set grows, its complement shrinks.
\end{remark*}

\begin{dependencies}
\begin{itemize}
  \item \hyperref[def:subset]{Subset}
  \item \hyperref[def:complement]{Complement}
\end{itemize}
\end{dependencies}

\begin{theorembox}{Theorem (Set Difference is Inclusion-Monotone in the Left Argument)}
\begin{theorem}[Set Difference is Inclusion-Monotone in the Left Argument]\label{thm:set-difference-monotone-left}
If $A\subseteq B$, then
\[
A\setminus C\subseteq B\setminus C.
\]
\hyperref[prf:set-difference-monotone-left]{Go to proof.}
\end{theorem}
\end{theorembox}

\LeanFormalizes{thm:set-difference-monotone-left}{lra-lean}{LRA.VolumeI.Set.Operations.Laws.Difference}{DifferenceMonotoneLeft}{pending}

\begin{remark*}[Standard quantified statement]
\[
  \forall A,B,C,\quad A\subseteq B\to A\setminus C\subseteq B\setminus C.
\]
\end{remark*}

\begin{remark*}[Predicate reading]
\[
  \operatorname{Subset}(A,B)
  \to
  \operatorname{Subset}(A\setminus C,B\setminus C).
\]
\end{remark*}

\begin{remark*}[Interpretation]
Set difference is inclusion-monotone in the set from which elements are removed.
\end{remark*}

\begin{dependencies}
\begin{itemize}
  \item \hyperref[def:subset]{Subset}
  \item \hyperref[def:set-difference]{Set Difference}
\end{itemize}
\end{dependencies}

\begin{theorembox}{Theorem (Set Difference is Inclusion-Antitone in the Right Argument)}
\begin{theorem}[Set Difference is Inclusion-Antitone in the Right Argument]\label{thm:set-difference-antitone-right}
If $C\subseteq D$, then
\[
A\setminus D\subseteq A\setminus C.
\]
\hyperref[prf:set-difference-antitone-right]{Go to proof.}
\end{theorem}
\end{theorembox}

\begin{remark*}[Standard quantified statement]
\[
  \forall A,C,D,\quad C\subseteq D\to A\setminus D\subseteq A\setminus C.
\]
\end{remark*}

\begin{remark*}[Predicate reading]
\[
  \operatorname{Subset}(C,D)
  \to
  \operatorname{Subset}(A\setminus D,A\setminus C).
\]
\end{remark*}

\begin{remark*}[Interpretation]
Set difference is inclusion-antitone in the set being removed.
\end{remark*}

\LeanFormalizes{thm:set-difference-antitone-right}{lra-lean}{LRA.VolumeI.Set.Operations.Laws.Difference}{DifferenceAntitoneRight}{pending}

\begin{dependencies}
\begin{itemize}
  \item \hyperref[def:subset]{Subset}
  \item \hyperref[def:set-difference]{Set Difference}
\end{itemize}
\end{dependencies}
